% ============================================================
% section11_final.tex - Раздел 11: Сквозной скрипт-оркестратор
% ============================================================
% !TeX program = xelatex
% Компиляция: xelatex --shell-escape main.tex

\section{Сквозной скрипт-оркестратор обработки данных} \label{sec:integrated_app}

В предыдущих разделах мы подробно изучили отдельные этапы обработки экспериментальных данных терагерцовой спектроскопии: от загрузки сырых сигналов до численного моделирования и пакетной оптимизации. В реальной работе исследователю крайне важно иметь возможность запустить сквозной цикл анализа, последовательно переходя от одного этапа к другому.

В данном разделе мы объединим все разработанные нами модули в единый, простой и эффективный скрипт-оркестратор \mycode{main.py}. Каждое графическое окно в этом конвейере запускается последовательно: после закрытия предыдущего окна управление автоматически передается следующей функции визуализации или расчета.

\subsection{Концепция последовательного запуска этапов (Pipeline)}

Поскольку все графические интерфейсы и визуализации в нашем проекте построены на стандартной библиотеке \mycode{matplotlib}, мы можем организовать сквозную обработку без использования сторонних GUI-фреймворков.

Принцип работы оркестратора предельно прост:
\begin{enumerate}
    \item Главная функция \mycode{main()} последовательно вызывает функции визуализации из других модулей (\mycode{plots.py}, \mycode{frequency_analysis.py}).
    \item Окно Matplotlib блокирует выполнение кода (вызов \mycode{plt.show()} приостанавливает поток).
    \item Пользователь интерактивно исследует графики (например, переключает углы с помощью радиокнопок).
    \item Как только пользователь закрывает текущее окно графика, выполнение возобновляется, и оркестратор немедленно запускает следующий шаг.
\end{enumerate}

Таким образом, главная функция \mycode{main()} превращается в простой и прозрачный конвейер:
\begin{minted}{python}
import numpy as np
from scipy.optimize import minimize
import config
import data_loader
import spectrum
import theoretical
import plots
import frequency_analysis

def main():
    print("=========================================================")
    print("ОРКЕСТРАТОР ОБРАБОТКИ ДАННЫХ ТГЦ-СПЕКТРОСКОПИИ (main.py)")
    print("=========================================================")
    
    # Шаг 1. Интерактивная фильтрация и спектры сигналов
    print("\n--- Шаг 1. Интерактивный анализ временных сигналов и спектров ---")
    print("Закройте окно графика, чтобы перейти к следующему шагу...")
    plots.main_interactive()
    
    # Шаг 2. Построение семейства спектров пропускания для всех углов
    print("\n--- Шаг 2. Визуализация семейства спектров пропускания ---")
    print("Закройте окно графика, чтобы перейти к следующему шагу...")
    plots.plot_all_transmissions()
    
    # Шаг 3. Локальная численная оптимизация на фиксированной частоте
    run_single_frequency_fit(target_freq=0.5)
    
    # Шаг 4. Теоретические характеристики поляризатора по модели Бланко
    print("\n--- Шаг 4. Теоретический расчет коэффициентов пропускания ---")
    print("Закройте окно графика, чтобы перейти к следующему шагу...")
    plots.plot_polarizer_theory_characteristics()
    
    # Шаг 5. Пакетная оптимизация по всему спектру частот
    print("\n--- Шаг 5. Пакетная оптимизация геометрических параметров по спектру ---")
    print("Запуск численного анализа. Пожалуйста, подождите...")
    results = frequency_analysis.analyze_polarizer_over_spectrum()
    print("Построение финального дашборда результатов...")
    frequency_analysis.plot_frequency_dependence(*results)
    
    print("\nВсе этапы обработки ТГц данных успешно завершены!")
\end{minted}

\subsubsection{Разбор кода: Главный конвейер main()}

Разберем структуру и последовательность вызовов в главной функции-оркестраторе:

\begin{syntaxdesc}
    \item[Интерактивный анализ \mycode{plots.main_interactive()}] Вызывает интерактивный дашборд для удаления постоянного смещения, настройки Гауссова временного фильтра и визуализации БПФ-спектров.
    \item[Семейство спектров \mycode{plots.plot_all_transmissions()}] Строит графики коэффициента пропускания по мощности для всех углов поворота поляризатора на одной координатной сетке.
    \item[Локальная оптимизация \mycode{run_single_frequency_fit}] Выполняет пятипараметрическую локальную подгонку модели Бланко под экспериментальные точки на частоте $0.5$ ТГц методом Нелдера-Мида.
    \item[Теоретические характеристики \mycode{plots.plot_polarizer_theory_characteristics()}] Строит теоретические амплитудные коэффициенты пропускания $|t_\perp|$ и $|t_\parallel|$ на логарифмической частотной шкале до $100$ ТГц, отмечая аномалию Вуда.
    \item[Пакетная оптимизация \mycode{frequency_analysis.analyze_polarizer_over_spectrum()}] Запускает пакетный расчет оптимальных параметров поляризатора по всему спектру с шагом по частоте.
\end{syntaxdesc}

\subsection{Описание этапов сквозной обработки}

Рассмотрим подробно каждый из этапов, реализуемых в нашем сквозном приложении.

\subsubsection{Шаг 1: Интерактивный анализ сигналов во временной и частотной областях}
Первый этап запускает интерактивный стенд (Рисунок~\ref{fig:9_1}), импортируемый из модуля \mycode{plots.py}. Пользователь может с помощью радиокнопок выбирать угол поворота ротатора, наглядно наблюдая исходный ТГц-импульс, форму наложенного Гауссова временного окна и спектры Фурье до и после фильтрации.

\smartfigure{fig9.1.png}{Интерактивная настройка параметров Гауссова временного фильтра на первом шаге приложения.}{fig:9_1}

\subsubsection{Шаг 2: Семейство спектров пропускания для всех углов}
На втором этапе строится семейство спектров пропускания для всех экспериментальных углов (Рисунок~\ref{fig:9_2}). Это позволяет исследователю качественно оценить динамический диапазон полезного терагерцового сигнала и выявить частотные резонансы (линии поглощения остаточного водяного пара в воздухе).

\smartfigure{fig9.2.png}{Визуализация семейства спектров пропускания для различных углов поворота поляризатора.}{fig:9_2}

\subsubsection{Шаг 3: Локальная подгонка параметров}
Скрипт выполняет точечную Nelder-Mead оптимизацию на частоте $0.5$ ТГц. На экран выводятся оптимальные шаг $P$ и диаметр проволок $D$, масштабный коэффициент периода и погрешность фиттинга. Это служит экспресс-проверкой сходимости численного алгоритма перед переходом к долгому пакетному расчету.

\subsubsection{Шаг 4: Теоретический анализ поляризатора}
Этот шаг строит теоретический спектральный отклик поляризатора (Рисунок~\ref{fig:9_3}) в широчайшем диапазоне частот (до $100$ ТГц). Студент наглядно видит область применимости квазистатического приближения (низкие частоты), а также положение резонансного провала (аномалии Вуда), где длина волны ТГц-излучения становится равной периоду решетки.

\smartfigure{fig9.3.png}{Теоретические спектры пропускания параллельной и перпендикулярной составляющих поляризатора в диапазоне 0.1--100 ТГц.}{fig:9_3}

\subsubsection{Шаг 5: Пакетный расчет по всему спектру частот}
На финальном этапе запускается пакетный оптимизатор (Рисунок~\ref{fig:9_4}). Он последовательно находит геометрические характеристики на каждой спектральной линии рабочего диапазона ($0.2$--$1.8$ ТГц) и отображает сводный пятипанельный дашборд.

\smartfigure{fig9.4.png}{Частотная зависимость оптимальных геометрических параметров проволочной решетки и коэффициента потерь.}{fig:9_4}

\begin{minitaskbox}
    \textbf{Исследуйте адекватность модели на краях диапазона} \\
    Запустите оркестратор \mycode{main.py}. Обратите внимание на финальные графики Шага 5.
    \begin{enumerate}
        \item Зафиксируйте поведение шага решетки $P(f)$ и диаметра $D(f)$ в диапазоне $0.4$--$1.4$ ТГц. Стабильны ли они?
        \item Посмотрите на поведение параметров ниже $0.2$ ТГц и выше $1.8$ ТГц. Почему там наблюдается хаотичный разброс точек?
        \item Объясните физический смысл роста коэффициента потерь с частотой.
    \end{enumerate}
\end{minitaskbox}

\subsection{Итоги раздела}
Созданный сквозной скрипт-оркестратор представляет собой законченное решение для автоматизированной экспресс-диагностики терагерцовых поляризаторов. Благодаря повторному использованию специализированных модулей (\mycode{data_loader.py}, \mycode{spectrum.py}, \mycode{plots.py}, \mycode{theoretical.py}), код оркестратора получился лаконичным, легко читаемым и расширяемым. Студент получает готовую исследовательскую платформу, на которой может проводить самостоятельную научную работу.
